% =====================================================================
%  Cours : CONTINUITÉ DES FONCTIONS
%  Terminale — Spécialité Mathématiques
%  Lycée Français International Jean CHARCOT — EL JADIDA
%  Année scolaire 2026 - 2027
%  Compilation : pdflatex (deux passes pour la table des matières)
% =====================================================================
\documentclass[11pt,a4paper]{article}

% ---------------------------------------------------------------------
%  Encodage, langue, polices
% ---------------------------------------------------------------------
\usepackage[utf8]{inputenc}
\usepackage[T1]{fontenc}
\usepackage{lmodern}
\usepackage[french]{babel}
\usepackage{microtype}

% ---------------------------------------------------------------------
%  Mathématiques
% ---------------------------------------------------------------------
\usepackage{amsmath,amssymb,mathtools}
\usepackage{mathrsfs}

% ---------------------------------------------------------------------
%  Mise en page
% ---------------------------------------------------------------------
\usepackage[a4paper,top=2.3cm,bottom=2.3cm,left=2cm,right=2cm,
            headheight=15pt]{geometry}
\usepackage[table]{xcolor}
\usepackage{graphicx}
\usepackage{array,booktabs,multicol}
\usepackage{enumitem}
\usepackage{fancyhdr}
\usepackage{titlesec}
\usepackage{lastpage}
\usepackage{needspace}
\usepackage{tocloft}
\setlength{\cftsecnumwidth}{2.6em}
\setlength{\cftsubsecnumwidth}{2em}
\setcounter{tocdepth}{2}
\setcounter{secnumdepth}{2}
\renewcommand{\cfttoctitlefont}{\hfill\Large\bfseries\color{bleu}}
\renewcommand{\cftaftertoctitle}{\hfill\hfill}
\renewcommand{\cftsecfont}{\bfseries\color{bleu}}
\renewcommand{\cftsecpagefont}{\bfseries\color{bleu}}
\tocloftpagestyle{fancy}
\usepackage{listings}
\usepackage[most]{tcolorbox}

% ---------------------------------------------------------------------
%  Figures
% ---------------------------------------------------------------------
\usepackage{tikz}
\usetikzlibrary{arrows.meta,calc,babel}
\usepackage{pgfplots}
\pgfplotsset{compat=1.18}
\usepackage{tkz-tab}

\usepackage[hidelinks]{hyperref}

% ---------------------------------------------------------------------
%  Couleurs
% ---------------------------------------------------------------------
\definecolor{bleu}{RGB}{20,55,120}
\definecolor{bleuclair}{RGB}{232,240,252}
\definecolor{rouge}{RGB}{170,30,45}
\definecolor{rougeclair}{RGB}{252,236,236}
\definecolor{vert}{RGB}{25,115,65}
\definecolor{vertclair}{RGB}{233,247,238}
\definecolor{orange}{RGB}{200,100,10}
\definecolor{orangeclair}{RGB}{254,243,228}
\definecolor{violet}{RGB}{95,45,140}
\definecolor{violetclair}{RGB}{243,236,250}
\definecolor{gris}{RGB}{90,90,90}
\definecolor{grisclair}{RGB}{244,244,244}
\definecolor{courbe}{RGB}{30,20,120}
\definecolor{magenta}{RGB}{170,30,140}

% ---------------------------------------------------------------------
%  Titres de sections
% ---------------------------------------------------------------------
\renewcommand{\thesection}{\Roman{section}}
\renewcommand{\thesubsection}{\arabic{subsection}.}
\newcommand{\boitesection}[1]{\colorbox{bleu}{\parbox{\dimexpr\linewidth-2\fboxsep}{%
  \strut\thesection~--~#1}}}
\newcommand{\boitesectionetoile}[1]{\colorbox{bleu}{\parbox{\dimexpr\linewidth-2\fboxsep}{%
  \strut#1}}}
\titleformat{\section}
  {\normalfont\Large\bfseries\color{white}}
  {}{0pt}{\boitesection}
\titleformat{name=\section,numberless}
  {\normalfont\Large\bfseries\color{white}}
  {}{0pt}{\boitesectionetoile}
\titlespacing*{\section}{0pt}{16pt}{10pt}
\titleformat{\subsection}
  {\normalfont\large\bfseries\color{bleu}}
  {\thesubsection}{0.6em}{}[{\color{bleu}\titlerule[0.6pt]}]
\titlespacing*{\subsection}{0pt}{12pt}{6pt}
\titleformat{\subsubsection}
  {\normalfont\normalsize\bfseries\color{bleu!80!black}}
  {}{0pt}{$\blacktriangleright$~}

% ---------------------------------------------------------------------
%  En-têtes et pieds de page
% ---------------------------------------------------------------------
\pagestyle{fancy}
\fancyhf{}
\fancyhead[L]{\small\color{gris}Cours -- Continuité des fonctions}
\fancyhead[R]{\small\color{gris}Terminale -- Spécialité Mathématiques}
\fancyfoot[L]{\small\color{gris}LFI Jean CHARCOT -- EL JADIDA}
\fancyfoot[C]{\small\color{gris}\thepage/\pageref{LastPage}}
\fancyfoot[R]{\small\color{gris}K. EL GHAZOULI}
\renewcommand{\headrulewidth}{0.4pt}
\renewcommand{\footrulewidth}{0.4pt}

% ---------------------------------------------------------------------
%  Environnements (tcolorbox)
% ---------------------------------------------------------------------
\tcbset{
  basecours/.style={enhanced, breakable, sharp corners=south,
    arc=2mm, boxrule=0.8pt, left=2mm, right=2mm, top=1.5mm, bottom=1.5mm,
    fonttitle=\bfseries, separator sign dash,
    before skip=8pt, after skip=8pt},
}

\newtcbtheorem[number within=section]{definition}{Définition}{
  basecours, colback=bleuclair, colframe=bleu, coltitle=white}{def}

\newtcbtheorem[use counter from=definition]{theoreme}{Théorème}{
  basecours, colback=rougeclair, colframe=rouge, coltitle=white}{thm}

\newtcbtheorem[use counter from=definition]{propriete}{Propriété}{
  basecours, colback=violetclair, colframe=violet, coltitle=white}{prop}

\newtcbtheorem[use counter from=definition]{corollaire}{Corollaire}{
  basecours, colback=rougeclair, colframe=rouge!80!black, coltitle=white}{cor}

\newtcbtheorem[number within=section]{exemple}{Exemple}{
  enhanced, breakable, colback=vertclair, colframe=vert, coltitle=vert,
  colbacktitle=vertclair, boxrule=0pt, leftrule=1.2mm, sharp corners,
  attach title to upper, after title={\ :\ }, fonttitle=\bfseries,
  separator sign dash, left=3mm, right=2mm, top=1.5mm, bottom=1.5mm,
  before skip=8pt, after skip=8pt}{ex}

\newtcolorbox{remarque}[1][Remarque]{
  enhanced, colback=grisclair, colframe=gris, boxrule=0pt,
  leftrule=1.2mm, sharp corners, left=3mm, right=2mm, top=1.5mm, bottom=1.5mm,
  title={#1}, attach title to upper, after title={\ :\ }, coltitle=gris,
  fonttitle=\bfseries, before skip=8pt, after skip=8pt}

\newtcolorbox{attention}{
  enhanced, breakable, colback=orangeclair, colframe=orange, boxrule=0.6pt,
  sharp corners, left=3mm, right=2mm, top=1.5mm, bottom=1.5mm,
  title={\large$\triangle$\ Attention}, attach title to upper,
  after title={\ :\ }, coltitle=orange, fonttitle=\bfseries,
  before skip=8pt, after skip=8pt}

\newtcolorbox{methode}[1][]{
  enhanced, breakable, colback=orangeclair, colframe=orange, coltitle=white,
  arc=2mm, boxrule=0.8pt, fonttitle=\bfseries, left=2mm, right=2mm,
  title={Méthode\ifx&#1&\else~--~#1\fi}, before skip=8pt, after skip=8pt}

\newtcolorbox{demo}{
  enhanced, breakable, colback=white, colframe=gris!60, boxrule=0.4pt,
  sharp corners, left=3mm, right=2mm, top=1.5mm, bottom=1.5mm,
  title={Démonstration}, attach title to upper, after title={\ :\ },
  coltitle=gris, fonttitle=\bfseries\itshape,
  before skip=8pt, after skip=8pt}

\newtcolorbox{programme}{
  enhanced, colback=bleuclair!50, colframe=bleu, boxrule=0.6pt, arc=2mm,
  title={Extrait du programme officiel (B.O.) -- Terminale spécialité},
  coltitle=white, colbacktitle=bleu, fonttitle=\bfseries,
  left=3mm, right=3mm}

% Fin de démonstration
\newcommand{\cqfd}{\hfill$\blacksquare$}

% ---------------------------------------------------------------------
%  Macros mathématiques
% ---------------------------------------------------------------------
\newcommand{\R}{\mathbb{R}}
\newcommand{\N}{\mathbb{N}}
\newcommand{\Cf}{\mathscr{C}_f}
\newcommand{\e}{\mathrm{e}}
\newcommand{\intff}[2]{\mathopen{[}#1\,;\,#2\mathclose{]}}
\newcommand{\intoo}[2]{\mathopen{]}#1\,;\,#2\mathclose{[}}
\newcommand{\intfo}[2]{\mathopen{[}#1\,;\,#2\mathclose{[}}
\newcommand{\intof}[2]{\mathopen{]}#1\,;\,#2\mathclose{]}}
\newcommand{\limite}[2]{\lim\limits_{#1}#2}
\newcommand{\suite}[1]{\left(#1_n\right)}

% ---------------------------------------------------------------------
%  Code Python
% ---------------------------------------------------------------------
\lstset{
  language=Python, basicstyle=\ttfamily\small, keywordstyle=\color{bleu}\bfseries,
  commentstyle=\color{vert}\itshape, stringstyle=\color{rouge},
  numbers=left, numberstyle=\tiny\color{gris}, frame=single,
  rulecolor=\color{gris!50}, backgroundcolor=\color{grisclair},
  showstringspaces=false, columns=fullflexible, keepspaces=true,
  literate={é}{{\'e}}1 {è}{{\`e}}1 {à}{{\`a}}1 {ê}{{\^e}}1
}

\setlist[itemize]{label=\textcolor{bleu}{$\bullet$}, leftmargin=*, itemsep=2pt}
\setlist[enumerate]{leftmargin=*, itemsep=2pt}
\setlength{\parindent}{0pt}
\setlength{\parskip}{4pt}

% =====================================================================
\begin{document}
% =====================================================================

% ---------------------------------------------------------------------
%  PAGE DE GARDE
% ---------------------------------------------------------------------
\begin{titlepage}
\begin{tikzpicture}[remember picture, overlay]
  \fill[bleu] (current page.north west) rectangle
              ([yshift=-0.9cm]current page.north east);
  \fill[bleu] (current page.south west) rectangle
              ([yshift=0.9cm]current page.south east);
  \draw[bleu, line width=1.2pt]
    ([xshift=1.2cm,yshift=-1.4cm]current page.north west) rectangle
    ([xshift=-1.2cm,yshift=1.4cm]current page.south east);
\end{tikzpicture}

\centering
\vspace*{0.6cm}

{\Large\bfseries\color{bleu} Lycée Français International Jean CHARCOT - EL JADIDA\par}
\vspace{0.5cm}
{\large Année scolaire 2026 - 2027\par}
\vspace{0.3cm}
{\large\bfseries Niveau : Terminale Spécialité Mathématiques\par}

\vspace{0.8cm}
{\color{bleu}\rule{0.85\linewidth}{1.2pt}}

\vspace{2.6cm}

\begin{tcolorbox}[enhanced, width=0.9\linewidth, colback=bleuclair,
  colframe=bleu, boxrule=1.5pt, arc=3mm, halign=center,
  top=8mm, bottom=8mm, drop shadow]
  {\fontsize{32}{40}\selectfont\bfseries\color{bleu}
   CONTINUITÉ\\[3mm] DES\hspace{0.4em}FONCTIONS\par}
  \vspace{4mm}
  {\large\itshape Cours, propriétés, démonstrations et exemples d'application\par}
\end{tcolorbox}

\vspace{1.4cm}

% Illustration : théorème des valeurs intermédiaires
\begin{tikzpicture}
\begin{axis}[axis lines=middle, width=10cm, height=6cm,
  xmin=0, xmax=6.2, ymin=0, ymax=3.6, xtick=\empty, ytick=\empty,
  xlabel={$x$}, ylabel={$y$}, clip=false,
  every axis x label/.style={at={(ticklabel* cs:1)}, anchor=west},
  every axis y label/.style={at={(ticklabel* cs:1)}, anchor=south}]
  \addplot[domain=1:5.6, samples=150, very thick, rouge]
     {2-sin(deg(pi*(x-2)/1.4))};
  \addplot[domain=0:6.2, thick, vert] {2};
  \foreach \c in {2,3.4,4.8}{
    \edef\temp{\noexpand\draw[dashed, gris] (axis cs:\c,0) -- (axis cs:\c,2);
               \noexpand\fill[bleu] (axis cs:\c,2) circle (2pt);}
    \temp
  }
  \node[left] at (axis cs:0,2) {$k$};
  \node[above, rouge] at (axis cs:5.6,1.05) {$\Cf$};
\end{axis}
\end{tikzpicture}

\vfill
{\large Enseignant : K. EL GHAZOULI\par}
\vspace{0.6cm}
\end{titlepage}

% ---------------------------------------------------------------------
%  SOMMAIRE ET PROGRAMME
% ---------------------------------------------------------------------
\setcounter{page}{1}
\begingroup
\hypersetup{linkcolor=bleu}
\renewcommand{\contentsname}{Sommaire}
\tableofcontents
\endgroup

\vspace{0.6cm}

\begin{programme}
\textbf{Contenus}
\begin{itemize}
  \item Fonction continue en un point (définition par la limite), sur un intervalle.
        Toute fonction dérivable est continue.
  \item Image d'une suite convergente par une fonction continue.
  \item Théorème des valeurs intermédiaires. Cas des fonctions continues
        strictement monotones.
\end{itemize}
\textbf{Capacités attendues}
\begin{itemize}
  \item Étudier les solutions d'une équation du type $f(x)=k$ :
        existence, unicité, encadrement.
  \item Pour une fonction continue $f$ d'un intervalle dans lui-même, étudier une
        suite définie par une relation de récurrence $u_{n+1}=f(u_n)$.
\end{itemize}
\textbf{Exemple d'algorithme} : méthode de dichotomie.\\
\textbf{Approfondissement} : exemples de fonctions discontinues (fonction partie
entière), de fonctions continues non dérivables.
\end{programme}

\begin{remarque}[Convention]
Dans tout ce chapitre, $I$ désigne un intervalle de $\R$ non vide et non réduit à
un point, et $\Cf$ la courbe représentative d'une fonction $f$ dans un repère
du plan.
\end{remarque}

% =====================================================================
\section{Continuité d'une fonction}
% =====================================================================

\subsection{Définition}

Soit $f$ une fonction définie sur un intervalle $I$ de $\R$. Intuitivement, dire
que $f$ est continue sur $I$ signifie que sa courbe représentative peut être
tracée \emph{en un seul morceau}, sans lever le crayon : la courbe ne présente
aucun saut, aucun trou. Mathématiquement, cela se traduit de la manière suivante.

\begin{definition}{Continuité en un point, sur un intervalle}{continuite}
Soit $f$ une fonction définie sur un intervalle $I$ de $\R$ et $a$ un réel de $I$.
\begin{itemize}
  \item On dit que $f$ est \textbf{continue en $a$} lorsque
  \[
    \limite{x\to a}{f(x)} = f(a).
  \]
  Lorsque $a$ n'est pas une borne de $I$, cela revient à :
  \[
    \limite{x\to a^-}{f(x)} \;=\; \limite{x\to a^+}{f(x)} \;=\; f(a).
  \]
  Sinon, on dit que $f$ est \textbf{discontinue} en $a$.
  \item On dit que $f$ est \textbf{continue sur $I$} lorsque $f$ est continue
  en \textbf{tout réel $a$ de $I$}.
\end{itemize}
\end{definition}

\begin{remarque}
Si $a$ est une borne de $I$ (par exemple $I=\intff{a}{b}$), on ne considère
que la limite à droite en $a$ : $f$ est continue en $a$ si
$\limite{x\to a^+}{f(x)}=f(a)$. De même en $b$ avec la limite à gauche.
\end{remarque}

\subsubsection{Illustration graphique}

Soit $f$ une fonction définie sur un intervalle $I$ de $\R$ et $a$ un réel de $I$.
On note $\Cf$ la courbe représentative de la fonction $f$ et $A$ le point de la
courbe $\Cf$ d'abscisse $a$. Pour tout réel $x$ de l'intervalle $I$, on considère
le point $M$ de la courbe $\Cf$ d'abscisse $x$.

\begin{center}
\begin{tabular}{@{}m{0.48\linewidth}@{\hspace{0.03\linewidth}}m{0.48\linewidth}@{}}
% ------------------ Figure 1 : fonction continue
\centering
\begin{tikzpicture}
\begin{axis}[axis lines=middle, width=7.6cm, height=6.2cm,
  xmin=-0.6, xmax=4.9, ymin=-0.8, ymax=3.6, xtick=\empty, ytick=\empty,
  xlabel={$x$}, ylabel={$y$}, clip=false,
  every axis x label/.style={at={(ticklabel* cs:1)}, anchor=west},
  every axis y label/.style={at={(ticklabel* cs:1)}, anchor=south}]
  \addplot[domain=-0.5:4.6, samples=200, thick, courbe]
     {x*(3.6-x)*exp(-0.35*(x-1.5)^2)};
  \draw[dashed, bleu] (axis cs:0,2.169) -- (axis cs:2.4,2.169) -- (axis cs:2.4,0);
  \draw[dashed, magenta] (axis cs:0,1.702) -- (axis cs:2.6,1.702) -- (axis cs:2.6,0);
  \fill[bleu] (axis cs:2.4,2.169) circle (1.6pt);
  \fill[magenta] (axis cs:2.6,1.702) circle (1.6pt);
  \node[left, bleu] at (axis cs:0,2.169) {$f(x)$};
  \node[left, magenta] at (axis cs:0,1.702) {$f(a)$};
  \node[below, bleu] at (axis cs:2.35,0) {$x$};
  \node[below, magenta] at (axis cs:2.7,0) {$a$};
  \node[above right, bleu] at (axis cs:2.4,2.169) {$M$};
  \node[above right, magenta] at (axis cs:2.6,1.702) {$A$};
  \node[below left] at (axis cs:0,0) {$O$};
  \node[below, courbe] at (axis cs:4.6,-0.2) {$\Cf$};
\end{axis}
\end{tikzpicture}
&
% ------------------ Figure 2 : fonction discontinue en a
\centering
\begin{tikzpicture}
\begin{axis}[axis lines=middle, width=7.6cm, height=6.2cm,
  xmin=-0.6, xmax=4.9, ymin=-0.8, ymax=3.6, xtick=\empty, ytick=\empty,
  xlabel={$x$}, ylabel={$y$}, clip=false,
  every axis x label/.style={at={(ticklabel* cs:1)}, anchor=west},
  every axis y label/.style={at={(ticklabel* cs:1)}, anchor=south}]
  \addplot[domain=-0.5:2.6, samples=200, thick, courbe]
     {x*(3.6-x)*exp(-0.35*(x-1.5)^2)};
  \addplot[domain=2.6:4.8, samples=60, thick, courbe] {0.5+0.25*(x-2.6)^2};
  \draw[thick, courbe, fill=white] (axis cs:2.6,1.702) circle (2pt);
  \draw[dashed, bleu] (axis cs:0,2.169) -- (axis cs:2.4,2.169) -- (axis cs:2.4,0);
  \draw[dashed, magenta] (axis cs:0,0.5) -- (axis cs:2.6,0.5) -- (axis cs:2.6,0);
  \fill[bleu] (axis cs:2.4,2.169) circle (1.6pt);
  \fill[magenta] (axis cs:2.6,0.5) circle (1.8pt);
  \node[left, bleu] at (axis cs:0,2.169) {$f(x)$};
  \node[left, magenta] at (axis cs:0,0.5) {$f(a)$};
  \node[below, bleu] at (axis cs:2.35,0) {$x$};
  \node[below, magenta] at (axis cs:2.7,0) {$a$};
  \node[above right, bleu] at (axis cs:2.4,2.169) {$M$};
  \node[above right, magenta] at (axis cs:2.6,0.5) {$A$};
  \node[below left] at (axis cs:0,0) {$O$};
  \node[above, courbe] at (axis cs:4.7,1.65) {$\Cf$};
\end{axis}
\end{tikzpicture}
\tabularnewline[2mm]
\textbf{La fonction $f$ est continue.}

Pour tout réel $a$ de $I$, on peut rendre $f(x)$ aussi proche que l'on veut
de $f(a)$ pourvu que $x$ soit suffisamment proche de $a$.
&
\textbf{La fonction $f$ n'est pas continue en $a$.}

La courbe $\Cf$ présente un saut au point d'abscisse $a$. Le point $M$ n'est pas
proche du point $A$ quand $x$ est proche de $a$.
\end{tabular}
\end{center}

\subsection{Exemples d'application}

\begin{exemple}{Continuité en un point par le calcul de limite}{}
Soit $f$ la fonction définie sur $\R$ par $f(x)=x^2-3x+1$. Montrons que $f$ est
continue en $2$.

D'une part, $f(2)=4-6+1=-1$. D'autre part, par somme et produit de limites,
$\limite{x\to 2}{f(x)} = 2^2-3\times 2+1=-1$.

Ainsi $\limite{x\to 2}{f(x)}=f(2)$ : la fonction $f$ est continue en $2$.
\end{exemple}

\begin{exemple}{Fonction définie par morceaux -- recherche d'un paramètre}{}
Soit $m$ un réel et $f$ la fonction définie sur $\R$ par
\[
  f(x)=\begin{cases}
    2x+1 & \text{si } x<1,\\[2pt]
    x^2+m & \text{si } x\geqslant 1.
  \end{cases}
\]
Sur $\intoo{-\infty}{1}$ et sur $\intoo{1}{+\infty}$, $f$ coïncide avec une
fonction polynôme, donc elle y est continue (voir paragraphe~\ref{sec:usuelles}).
Étudions la continuité en $1$ :
\begin{itemize}
  \item $\limite{x\to 1^-}{f(x)}=\limite{x\to 1^-}{(2x+1)}=3$ ;
  \item $\limite{x\to 1^+}{f(x)}=\limite{x\to 1^+}{(x^2+m)}=1+m=f(1)$.
\end{itemize}
Donc $f$ est continue en $1$ si et seulement si $1+m=3$, c'est-à-dire $m=2$.

\textbf{Conclusion :} $f$ est continue sur $\R$ si et seulement si $m=2$.
\end{exemple}

\begin{exemple}{Une fonction discontinue : la fonction partie entière}{}
Pour tout réel $x$, il existe un unique entier relatif $n$ tel que
$n\leqslant x<n+1$ ; on l'appelle \textbf{partie entière} de $x$ et on le note
$E(x)$ (en Python : \texttt{floor(x)}). Par exemple $E(2{,}7)=2$, $E(3)=3$,
$E(-1{,}2)=-2$.

\begin{minipage}{0.52\linewidth}
Étudions la continuité en $1$ :
\begin{itemize}
  \item pour $x\in\intfo{0}{1}$, $E(x)=0$, donc $\limite{x\to 1^-}{E(x)}=0$ ;
  \item pour $x\in\intfo{1}{2}$, $E(x)=1$, donc $\limite{x\to 1^+}{E(x)}=1$.
\end{itemize}
Les limites à gauche et à droite en $1$ sont différentes : la fonction $E$
\textbf{n'est pas continue en $1$}.

Plus généralement, $E$ est discontinue en tout entier relatif $n$, mais elle est
continue sur chaque intervalle $\intfo{n}{n+1}$ (où elle est constante).
La courbe présente un « saut » en chaque entier.
\end{minipage}\hfill
\begin{minipage}{0.44\linewidth}
\centering
\begin{tikzpicture}[scale=0.62]
  \draw[gris!30] (-2.5,-2.5) grid (3.5,3.5);
  \draw[-{Stealth}] (-2.7,0) -- (3.8,0) node[right] {$x$};
  \draw[-{Stealth}] (0,-2.7) -- (0,3.8) node[above] {$y$};
  \foreach \k in {-2,-1,1,2,3}{
    \node[below, font=\scriptsize] at (\k,0) {$\k$};
    \node[left,  font=\scriptsize] at (0,\k) {$\k$};}
  \node[below left, font=\scriptsize] at (0,0) {$O$};
  \foreach \n in {-2,-1,0,1,2}{
    \draw[very thick, rouge] (\n,\n) -- (\n+1,\n);
    \fill[rouge] (\n,\n) circle (2.6pt);
    \draw[thick, rouge, fill=white] (\n+1,\n) circle (2.6pt);}
  \node[rouge, above] at (2.2,3.4) {$y=E(x)$};
\end{tikzpicture}
\end{minipage}
\end{exemple}

% =====================================================================
\section{Continuité des fonctions usuelles et opérations}\label{sec:usuelles}
% =====================================================================

\needspace{9cm}
\subsection{Continuité des fonctions de référence}

On a le tableau de continuité des fonctions de référence suivant (résultats admis) :

\begin{center}
\renewcommand{\arraystretch}{1.55}
\begin{tabular}{|>{\centering\arraybackslash}m{6.6cm}|>{\centering\arraybackslash}m{6.6cm}|}
\hline
\rowcolor{bleuclair}\textbf{Fonctions usuelles} & \textbf{Intervalle(s) de continuité} \\ \hline
$f(x)=x^n$, \quad $n\in\N$ & $\R$ \\ \hline
$f(x)=\dfrac{1}{x^n}$, \quad $n\in\N^*$ & $\intoo{-\infty}{0}$ et $\intoo{0}{+\infty}$ \\ \hline
$f(x)=\sqrt{x}$ & $\intfo{0}{+\infty}$ \\ \hline
$f(x)=\lvert x\rvert$ & $\R$ \\ \hline
$f(x)=\e^x$ & $\R$ \\ \hline
$f(x)=\sin x$ \ et \ $f(x)=\cos x$ & $\R$ \\ \hline
$f(x)=\ln x$ \ {\small(chapitre « Logarithme »)} & $\intoo{0}{+\infty}$ \\ \hline
\end{tabular}
\end{center}

\begin{remarque}
Il en découle (avec le théorème suivant) que toute \textbf{fonction polynôme}
est continue sur $\R$ et que toute \textbf{fonction rationnelle} (quotient de
deux polynômes) est continue sur chacun des intervalles où elle est définie.
\end{remarque}

\subsection{Opérations sur les fonctions continues}

\begin{theoreme}{Opérations}{operations}
Soit $k$ un réel. Soit $f$ et $g$ deux fonctions définies sur un intervalle $I$
et soit $a\in I$. Si $f$ et $g$ sont continues en $a$, alors :
\begin{itemize}
  \item $f+g$ est continue en $a$ ;
  \item $k\times f$ est continue en $a$ ;
  \item $f\times g$ est continue en $a$ ;
  \item $\dfrac{f}{g}$ est continue en $a$, à condition que $g(a)\neq 0$.
\end{itemize}
\textbf{Composition :} si $f$ est continue en $a$ et si $g$ est définie sur un
intervalle $J$ contenant $f(I)$ et continue en $f(a)$, alors $g\circ f$ est
continue en $a$.

Ces résultats s'étendent aux fonctions continues sur un intervalle.
\end{theoreme}

\subsection{Exemples d'application}

\begin{exemple}{Fonction rationnelle}{}
La fonction $f:x\longmapsto \dfrac{3x^2-1}{x-2}$ est le quotient de deux fonctions
polynômes, continues sur $\R$, et le dénominateur ne s'annule qu'en $2$.
Donc $f$ est continue sur $\intoo{-\infty}{2}$ et sur $\intoo{2}{+\infty}$.
\end{exemple}

\begin{exemple}{Composition}{}
Soit $g$ la fonction définie sur $\R$ par $g(x)=\sqrt{x^2+1}$.

La fonction $u:x\longmapsto x^2+1$ est continue sur $\R$ (polynôme) et, pour tout
réel $x$, $u(x)\geqslant 1>0$ : $u(\R)\subset\intfo{0}{+\infty}$.
La fonction racine carrée est continue sur $\intfo{0}{+\infty}$.
Par composition, $g=\sqrt{\phantom{x}}\circ u$ est continue sur $\R$.
\end{exemple}

\begin{exemple}{Produit et composition}{}
Soit $h$ la fonction définie sur $\R$ par $h(x)=\e^{-x}\cos(2x)$.
\begin{itemize}
  \item $x\mapsto -x$ est continue sur $\R$ et $\exp$ est continue sur $\R$,
        donc $x\mapsto \e^{-x}$ est continue sur $\R$ (composition) ;
  \item $x\mapsto 2x$ et $\cos$ sont continues sur $\R$, donc
        $x\mapsto\cos(2x)$ est continue sur $\R$ (composition).
\end{itemize}
Par produit, $h$ est continue sur $\R$.
\end{exemple}

% =====================================================================
\section{Continuité et dérivabilité}
% =====================================================================

\subsection{Théorème}

\begin{theoreme}{Dérivabilité implique continuité}{derivcont}
Soit $f$ une fonction définie sur un intervalle $I$ et soit $a\in I$.
\begin{itemize}
  \item Si $f$ est \textbf{dérivable en $a$}, alors la fonction $f$ est
        \textbf{continue en $a$}.
  \item Si $f$ est \textbf{dérivable sur $I$}, alors la fonction $f$ est
        \textbf{continue sur $I$}.
\end{itemize}
\end{theoreme}

\begin{demo}
Supposons $f$ dérivable en $a$. Pour tout réel $h\neq 0$ tel que $a+h\in I$, on a
\[
  f(a+h) = f(a) + h\times\frac{f(a+h)-f(a)}{h}.
\]
Or $\limite{h\to 0}{\dfrac{f(a+h)-f(a)}{h}}=f'(a)$ (réel) et
$\limite{h\to 0}{h}=0$, donc par produit
$\limite{h\to 0}{h\times\dfrac{f(a+h)-f(a)}{h}}=0\times f'(a)=0$.
Ainsi $\limite{h\to 0}{f(a+h)}=f(a)$, c'est-à-dire $\limite{x\to a}{f(x)}=f(a)$ :
$f$ est continue en $a$. \cqfd
\end{demo}

\subsection{La réciproque est fausse}

\begin{minipage}{0.58\linewidth}
\begin{remarque}
La réciproque de ce théorème est fausse : une fonction peut être continue en un
point $a$ sans pour autant y être dérivable.

C'est le cas de la fonction valeur absolue : $f(x)=\lvert x\rvert$.

En effet, $f$ est continue en $0$ (car $\limite{x\to 0}{\lvert x\rvert}=0=f(0)$)
mais elle n'est pas dérivable en $0$ : pour $h\neq 0$,
\[
  \frac{f(0+h)-f(0)}{h}=\frac{\lvert h\rvert}{h}=
  \begin{cases} 1 & \text{si } h>0,\\ -1 & \text{si } h<0.\end{cases}
\]
Le taux de variation admet deux limites différentes à droite et à gauche
en $0$ ($1$ et $-1$) : il n'a donc pas de limite en $0$.
\end{remarque}
\end{minipage}\hfill
\begin{minipage}{0.38\linewidth}
\centering
\begin{tikzpicture}[scale=0.62]
  \draw[gris!30] (-3.5,-0.8) grid (3.5,3.8);
  \draw[-{Stealth}] (-3.7,0) -- (3.8,0);
  \draw[-{Stealth}] (0,-0.9) -- (0,4);
  \foreach \k in {-3,-2,-1,1,2,3}{\node[below, font=\scriptsize] at (\k,0) {$\k$};}
  \foreach \k in {1,2,3}{\node[left, font=\scriptsize] at (0,\k) {$\k$};}
  \node[below left, font=\scriptsize] at (0,0) {$O$};
  \draw[thick, rouge] (-3.6,3.6) -- (0,0) -- (3.6,3.6);
  \draw[-{Stealth}, very thick, bleu] (0,0) -- (-1,1);
  \draw[-{Stealth}, very thick, bleu] (0,0) -- (1,1);
  \node[bleu, left, font=\small, fill=white, inner sep=1pt]  at (-1.05,1.05) {pente $-1$};
  \node[bleu, right, font=\small, fill=white, inner sep=1pt] at (1.05,1.05) {pente $+1$};
\end{tikzpicture}

\small La courbe est en « un seul morceau » \\ et possède un point anguleux en $O$.
\end{minipage}

\subsection{Exemples d'application}

\begin{exemple}{Utiliser la dérivabilité pour justifier la continuité}{}
La fonction $f$ définie sur $\R$ par $f(x)=(2x-1)\e^{x}$ est dérivable sur $\R$
comme produit de fonctions dérivables sur $\R$. D'après le théorème~\ref{thm:derivcont},
$f$ est donc continue sur $\R$.
\end{exemple}

\begin{exemple}{La fonction racine carrée en $0$}{}
La fonction $f:x\mapsto\sqrt{x}$ est continue en $0$ car
$\limite{x\to 0^+}{\sqrt{x}}=0=f(0)$. En revanche, pour $h>0$ :
\[
  \frac{f(0+h)-f(0)}{h}=\frac{\sqrt{h}}{h}=\frac{1}{\sqrt{h}}
  \xrightarrow[h\to 0^+]{}+\infty .
\]
Le taux de variation n'a pas de limite finie : $f$ n'est pas dérivable en $0$.
Graphiquement, la courbe admet en $O$ une tangente verticale.
\end{exemple}

% =====================================================================
\section{Continuité et suites}
% =====================================================================

\subsection{Image d'une suite convergente par une fonction continue}

\begin{theoreme}{Image d'une suite convergente}{imagesuite}
Soit $\suite{u}$ une suite à valeurs dans un intervalle $I$, convergeant vers un
réel $\ell$ de $I$, et $f$ une fonction définie sur $I$.

Si $f$ est \textbf{continue en $\ell$}, alors la suite $\big(f(u_n)\big)$ converge
vers $f(\ell)$.
\end{theoreme}

\begin{exemple}{}{}
Soit $\suite{v}$ la suite définie sur $\N$ par $v_n=\sqrt{4+\dfrac{1}{n+1}}$.

On a $v_n=f(u_n)$ avec $u_n=4+\dfrac{1}{n+1}$ et $f(x)=\sqrt{x}$.
Or $\limite{n\to+\infty}{u_n}=4$ et $f$ est continue en $4$, donc
$\limite{n\to+\infty}{v_n}=\sqrt{4}=2$.
\end{exemple}

\begin{attention}
L'hypothèse de continuité est indispensable. Avec la fonction partie entière $E$,
discontinue en $1$, et $u_n=1-\dfrac{1}{n+1}$ (pour $n\geqslant 1$, $0<u_n<1$) :
$\limite{n\to+\infty}{u_n}=1$ mais $E(u_n)=0$ pour tout $n\geqslant 1$, donc
$\big(E(u_n)\big)$ converge vers $0\neq E(1)=1$.
\end{attention}

\subsection{Théorème du point fixe}

\begin{theoreme}{Théorème du point fixe}{pointfixe}
Soit une suite $\suite{u}$ définie par $u_0\in I$ et la relation
$u_{n+1}=f(u_n)$, où $f$ est une fonction telle que $f(I)\subset I$,
convergente vers $\ell\in I$.

Si la fonction associée $f$ est continue en $\ell$, alors la limite $\ell$ de la
suite $\suite{u}$ est solution de l'équation $f(x)=x$.
\end{theoreme}

\begin{demo}
D'une part, $\limite{n\to+\infty}{u_{n+1}}=\ell$. D'autre part, $f$ étant
continue en $\ell$, le théorème~\ref{thm:imagesuite} donne
$\limite{n\to+\infty}{f(u_n)}=f(\ell)$. Comme $u_{n+1}=f(u_n)$ pour tout $n$,
l'unicité de la limite donne $f(\ell)=\ell$. \cqfd
\end{demo}

\begin{remarque}
On dit que $\ell$ est un \textbf{point fixe} de $f$. Ce théorème ne prouve pas
la convergence : il permet seulement, \emph{une fois la convergence établie},
de déterminer les valeurs possibles de la limite.
\end{remarque}

\subsection{Exemples d'application}

\begin{exemple}{Suite définie par une racine}{}
Soit une suite $\suite{u}$ définie sur $\N$ par :
$\begin{cases} u_0=0\\ u_{n+1}=\sqrt{3u_n+4}\end{cases}$

On admet que $\suite{u}$ est positive, croissante et majorée par $4$ : pour tout
$n\in\N$, $0\leqslant u_n\leqslant 4$. D'après le théorème de convergence des
suites monotones, $\suite{u}$ est convergente vers une limite $\ell$, et par
passage à la limite dans les inégalités, $0\leqslant\ell\leqslant 4$.

\textbf{Déterminons $\ell$.}

On considère la fonction $f$ définie sur $\intfo{-\frac{4}{3}}{+\infty}$ par
$f(x)=\sqrt{3x+4}$. La fonction $f$ est continue sur $\intfo{-\frac{4}{3}}{+\infty}$
(composée de $x\mapsto 3x+4$ et de la racine carrée), donc sur $\intff{0}{4}$ ;
de plus $f$ est croissante sur $\intff{0}{4}$ avec $f(0)=2$ et $f(4)=4$, donc
$f\big(\intff{0}{4}\big)\subset\intff{0}{4}$.

D'après le théorème du point fixe, $\ell$ est solution de l'équation $f(x)=x$.
Pour $x\in\intff{0}{4}$ :
\[
  \begin{aligned}
  \sqrt{3x+4}=x \;&\Longrightarrow\; 3x+4=x^2
  \iff x^2-3x-4=0\\
  &\iff (x+1)(x-4)=0
  \iff x=-1 \ \text{ ou } \ x=4 .
  \end{aligned}
\]
Comme $u_n\geqslant 0$ pour tout $n$, on a $\ell\geqslant 0$ et la seule valeur
possible est $4$ ; on vérifie que $\sqrt{3\times 4+4}=\sqrt{16}=4$.

\textbf{Conclusion :} la suite $\suite{u}$ converge vers $4$.
\end{exemple}

\begin{exemple}{Algorithme de Héron : approximation de $\sqrt{2}$}{}
Soit la suite $\suite{u}$ définie par $u_0=2$ et, pour tout $n\in\N$,
$u_{n+1}=\dfrac{1}{2}\left(u_n+\dfrac{2}{u_n}\right)$.

On admet que, pour tout $n\in\N$, $u_n\geqslant 1$ et que $\suite{u}$ est
décroissante. Étant décroissante et minorée par $1$, elle converge vers un réel
$\ell\geqslant 1$.

La fonction $f:x\mapsto \dfrac{1}{2}\left(x+\dfrac{2}{x}\right)$ est continue sur
$\intfo{1}{+\infty}$ (fonction rationnelle définie sur cet intervalle), donc en
$\ell$. D'après le théorème du point fixe :
\[
  f(\ell)=\ell \iff \frac{1}{2}\left(\ell+\frac{2}{\ell}\right)=\ell
  \iff \frac{2}{\ell}=\ell \iff \ell^2=2 .
\]
Comme $\ell\geqslant 1$, on obtient $\ell=\sqrt{2}$. Les premiers termes
$u_1=1{,}5$ ; $u_2\approx 1{,}416\,67$ ; $u_3\approx 1{,}414\,216$ fournissent
d'excellentes approximations de $\sqrt{2}$.
\end{exemple}

% =====================================================================
\section{Continuité et résolution d'équations}
% =====================================================================

\subsection{Théorème des valeurs intermédiaires}

\begin{theoreme}{Théorème des valeurs intermédiaires (TVI) -- cas général}{tvi}
Soit $f$ une fonction \textbf{continue} sur un intervalle $\intff{a}{b}$ et $k$
un réel compris entre $f(a)$ et $f(b)$.

Alors il existe (au moins) un réel $c$ dans $\intff{a}{b}$ tel que $f(c)=k$.

Autrement dit, l'équation $f(x)=k$ admet au moins une solution dans $\intff{a}{b}$.
\end{theoreme}

\begin{minipage}{0.55\linewidth}
\begin{remarque}
Ce théorème (admis) résulte du fait que \textbf{l'image d'un intervalle} de $\R$
par une fonction continue est \textbf{un intervalle} de $\R$.
\end{remarque}

\textbf{Illustration graphique :}

Ci-contre, $k$ est bien compris entre $f(a)$ et $f(b)$. L'équation $f(x)=k$
admet donc des solutions.

\textbf{L'existence de $c$ ne veut pas dire qu'il soit unique.}
Ici, il existe 3 valeurs : $c_1$, $c_2$ et $c_3$.
\end{minipage}\hfill
\begin{minipage}{0.42\linewidth}
\centering
\begin{tikzpicture}
\begin{axis}[axis lines=middle, width=7.2cm, height=5.6cm,
  xmin=0, xmax=6.2, ymin=0, ymax=3.5, xtick=\empty, ytick=\empty, clip=false,
  axis line style={-{Stealth}}]
  \draw[line width=3pt, rouge] (axis cs:0,1.025) -- (axis cs:0,2.782);
  \draw[line width=3pt, bleu] (axis cs:1,0) -- (axis cs:5.6,0);
  \addplot[domain=0:6.2, thick, vert] {2};
  \addplot[domain=1:5.6, samples=150, thick, black]
     {2-sin(deg(pi*(x-2)/1.4))};
  \draw[dashed] (axis cs:0,2.782) -- (axis cs:1,2.782) -- (axis cs:1,0);
  \draw[dashed] (axis cs:0,1.025) -- (axis cs:5.6,1.025) -- (axis cs:5.6,0);
  \foreach \c in {2,3.4,4.8}{
    \edef\temp{\noexpand\draw[dashed] (axis cs:\c,0) -- (axis cs:\c,2);}
    \temp}
  \node[left] at (axis cs:0,2.782) {$f(a)$};
  \node[left] at (axis cs:0,2) {$k$};
  \node[left] at (axis cs:0,1.025) {$f(b)$};
  \node[below] at (axis cs:1,0) {$a$};
  \node[below] at (axis cs:2,0) {$c_1$};
  \node[below] at (axis cs:3.4,0) {$c_2$};
  \node[below] at (axis cs:4.8,0) {$c_3$};
  \node[below] at (axis cs:5.6,0) {$b$};
  \node[below left] at (axis cs:0,0) {$O$};
\end{axis}
\end{tikzpicture}
\end{minipage}

\begin{exemple}{Existence d'une solution}{}
Montrons que l'équation $\cos x = x$ admet au moins une solution dans
$\intff{0}{\frac{\pi}{2}}$.

Posons $h(x)=\cos x - x$. La fonction $h$ est continue sur
$\intff{0}{\frac{\pi}{2}}$ (différence de fonctions continues) et :
\[
  h(0)=1>0 \qquad\text{et}\qquad h\!\left(\tfrac{\pi}{2}\right)=0-\tfrac{\pi}{2}<0 .
\]
Ainsi $0$ est compris entre $h\big(\tfrac{\pi}{2}\big)$ et $h(0)$. D'après le TVI,
il existe au moins un réel $c\in\intff{0}{\frac{\pi}{2}}$ tel que $h(c)=0$,
c'est-à-dire $\cos c=c$.
\end{exemple}

\subsection{Extension aux intervalles ouverts}

Il est également possible d'utiliser les limites dans le théorème des valeurs
intermédiaires.

\begin{theoreme}{TVI avec des limites}{tvilim}
Soit $f$ une fonction continue sur un intervalle $\intoo{a}{b}$ (avec $a$ et $b$
réels ou infinis) telle que $\limite{x\to a^+}{f(x)}$ et $\limite{x\to b^-}{f(x)}$
existent (finies ou infinies).

Soit $k$ un réel strictement compris entre ces deux limites.

Alors il existe (au moins) un réel $c$ dans $\intoo{a}{b}$ tel que $f(c)=k$.
\end{theoreme}

\begin{exemple}{}{}
Soit $f$ la fonction polynôme de degré $3$ définie sur $\R$ par $f(x)=x^3+3$.

$f$ est continue sur $\R$ (polynôme), et
$\limite{x\to+\infty}{f(x)}=+\infty$, \ $\limite{x\to-\infty}{f(x)}=-\infty$.

Or $0\in\intoo{-\infty}{+\infty}$. D'après le théorème des valeurs intermédiaires,
il existe un réel $x$ tel que $f(x)=0$ (ici $x=-\sqrt[3]{3}$).
\end{exemple}

\subsection{Cas des fonctions strictement monotones}

\begin{corollaire}{Corollaire du TVI -- théorème de la bijection}{bijection}
Soit $f$ une fonction \textbf{continue} et \textbf{strictement monotone} sur
$I=\intff{a}{b}$.

Pour tout réel $k$ compris entre $f(a)$ et $f(b)$, l'équation $f(x)=k$ admet
\textbf{une unique solution} dans $I$.
\end{corollaire}

\begin{minipage}{0.56\linewidth}
\begin{demo}
\emph{Existence :} c'est le TVI. \emph{Unicité :} si $c\neq c'$ étaient deux
solutions, par exemple $c<c'$, la stricte monotonie donnerait $f(c)<f(c')$
ou $f(c)>f(c')$, ce qui contredit $f(c)=f(c')=k$. \cqfd
\end{demo}

\begin{remarque}
\begin{itemize}
  \item Ce théorème est aussi appelé \textbf{théorème de la bijection}.
  \item On généralise ce théorème à l'intervalle ouvert $I=\intoo{a}{b}$.
        Le réel $k$ doit alors être strictement compris entre
        $\limite{x\to a^+}{f(x)}$ et $\limite{x\to b^-}{f(x)}$.
  \item Lorsque $k=0$, on pourra montrer que $f(a)\times f(b)<0$.
\end{itemize}
\end{remarque}
\end{minipage}\hfill
\begin{minipage}{0.40\linewidth}
\centering
\begin{tikzpicture}
\begin{axis}[axis lines=middle, width=7cm, height=5.8cm,
  xmin=0, xmax=6.3, ymin=0, ymax=4.9, xtick=\empty, ytick=\empty, clip=false,
  axis line style={-{Stealth}}]
  \draw[line width=3pt, rouge] (axis cs:0,0.92) -- (axis cs:0,4.43);
  \draw[line width=3pt, bleu] (axis cs:1,0) -- (axis cs:5.5,0);
  \addplot[domain=0:6.3, thick, vert] {3};
  \addplot[domain=1:5.5, samples=100, thick, rouge] {0.8+0.12*x^2};
  \draw[dashed] (axis cs:0,0.92) -- (axis cs:1,0.92) -- (axis cs:1,0);
  \draw[dashed] (axis cs:0,4.43) -- (axis cs:5.5,4.43) -- (axis cs:5.5,0);
  \draw[dashed] (axis cs:4.282,3) -- (axis cs:4.282,0);
  \fill (axis cs:4.282,3) circle (1.6pt);
  \node[left] at (axis cs:0,4.43) {$f(b)$};
  \node[left] at (axis cs:0,3) {$k$};
  \node[left] at (axis cs:0,0.92) {$f(a)$};
  \node[below] at (axis cs:1,0) {$a$};
  \node[below] at (axis cs:4.282,0) {$c$};
  \node[below] at (axis cs:5.5,0) {$b$};
  \node[below left] at (axis cs:0,0) {$O$};
\end{axis}
\end{tikzpicture}
\end{minipage}

\begin{remarque}[Convention des tableaux de variations]
Dans un tableau de variations, une flèche oblique traduit la
\textbf{continuité} et la \textbf{stricte monotonie} de la fonction sur
l'intervalle correspondant.
\end{remarque}

\begin{methode}[Étudier une équation $f(x)=k$ : existence, unicité, encadrement]
\begin{enumerate}
  \item Justifier que $f$ est \textbf{continue} sur l'intervalle $I$ considéré.
  \item Justifier que $f$ est \textbf{strictement monotone} sur $I$
        (souvent par l'étude du signe de $f'$).
  \item Calculer les images (ou les limites) aux bornes de $I$ et vérifier que
        $k$ est compris entre ces valeurs.
  \item Conclure avec le corollaire du TVI : l'équation $f(x)=k$ admet une
        unique solution $\alpha$ dans $I$.
  \item \textbf{Encadrer} $\alpha$ à l'aide de la calculatrice (balayage,
        tableau de valeurs) ou d'un algorithme (dichotomie).
\end{enumerate}
\end{methode}

\subsection{Exemples d'application}

\begin{exemple}{Existence, unicité et encadrement d'une solution}{}
Soit $f$ la fonction définie sur $\R$ par $f(x)=x^3+x-1$.
Montrons que l'équation $f(x)=0$ admet une unique solution $\alpha$ dans $\R$,
puis donnons-en un encadrement d'amplitude $10^{-2}$.

\begin{itemize}
  \item $f$ est une fonction polynôme, donc continue et dérivable sur $\R$.
  \item Pour tout réel $x$, $f'(x)=3x^2+1>0$ : $f$ est strictement croissante sur $\R$.
  \item $\limite{x\to-\infty}{f(x)}=\limite{x\to-\infty}{x^3}=-\infty$ et
        $\limite{x\to+\infty}{f(x)}=\limite{x\to+\infty}{x^3}=+\infty$.
\end{itemize}

\begin{center}
\begin{tikzpicture}
\tkzTabInit[lgt=2.2, espcl=4]{$x$/1, $f'(x)$/1, $f(x)$/2}{$-\infty$, $+\infty$}
\tkzTabLine{, +, }
\tkzTabVar{-/$-\infty$, +/$+\infty$}
\tkzTabVal{1}{2}{0.5}{$\alpha$}{$0$}
\end{tikzpicture}
\end{center}

Comme $0\in\intoo{-\infty}{+\infty}$, d'après le corollaire du TVI, l'équation
$f(x)=0$ admet une unique solution $\alpha$ dans $\R$.

\textbf{Encadrement :} $f(0)=-1<0$ et $f(1)=1>0$, donc $0<\alpha<1$.
Par balayage avec un pas de $0{,}01$ :
\[
  f(0{,}68)\approx -0{,}0056<0 \qquad\text{et}\qquad f(0{,}69)\approx 0{,}0185>0,
\]
donc, $f$ étant strictement croissante : $\boxed{0{,}68<\alpha<0{,}69}$.
\end{exemple}

\begin{exemple}{Dénombrer les solutions à l'aide du tableau de variations}{}
Soit $g$ la fonction définie sur $\R$ par $g(x)=x^3-3x+1$. Déterminons le nombre
de solutions de l'équation $g(x)=0$.

$g$ est continue et dérivable sur $\R$ et $g'(x)=3x^2-3=3(x-1)(x+1)$.
De plus $g(-1)=3$, $g(1)=-1$,
$\limite{x\to-\infty}{g(x)}=-\infty$ et $\limite{x\to+\infty}{g(x)}=+\infty$.

\begin{center}
\begin{tikzpicture}
\tkzTabInit[lgt=2.2, espcl=2.8]{$x$/1, $g'(x)$/1, $g(x)$/2.2}
  {$-\infty$, $-1$, $1$, $+\infty$}
\tkzTabLine{, +, z, -, z, +, }
\tkzTabVar{-/$-\infty$, +/$3$, -/$-1$, +/$+\infty$}
\end{tikzpicture}
\end{center}

\begin{itemize}
  \item Sur $\intof{-\infty}{-1}$ : $g$ est continue, strictement croissante,
        et $0\in\intof{-\infty}{3}$ ; d'où une unique solution $\alpha_1$.
  \item Sur $\intff{-1}{1}$ : $g$ est continue, strictement décroissante, et
        $0\in\intff{-1}{3}$ ; d'où une unique solution $\alpha_2$.
  \item Sur $\intfo{1}{+\infty}$ : $g$ est continue, strictement croissante, et
        $0\in\intfo{-1}{+\infty}$ ; d'où une unique solution $\alpha_3$.
\end{itemize}
L'équation $g(x)=0$ admet donc exactement \textbf{trois solutions} dans $\R$.
Comme $g(-2)=-1$, $g(0)=1$ et $g(2)=3$, on a de plus :
$-2<\alpha_1<-1$, \ $0<\alpha_2<1$ \ et \ $1<\alpha_3<2$.
\end{exemple}

\subsection{Algorithme de dichotomie}

On suppose que $f$ est continue et strictement monotone sur $\intff{a}{b}$, avec
$f(a)\times f(b)<0$ ; l'équation $f(x)=0$ admet donc une unique solution
$\alpha\in\intff{a}{b}$. Le principe de la \textbf{dichotomie} consiste à couper
l'intervalle en deux à chaque étape :
\begin{itemize}
  \item on calcule le milieu $m=\dfrac{a+b}{2}$ ;
  \item si $f(a)\times f(m)\leqslant 0$, alors $\alpha\in\intff{a}{m}$ et on
        remplace $b$ par $m$ ; sinon $\alpha\in\intff{m}{b}$ et on remplace $a$
        par $m$ ;
  \item on recommence tant que l'amplitude $b-a$ est supérieure à la précision
        $p$ souhaitée.
\end{itemize}
Après $n$ étapes, l'amplitude de l'intervalle est $\dfrac{b-a}{2^n}$.

\begin{minipage}[t]{0.55\linewidth}
\begin{lstlisting}
def dichotomie(f, a, b, p):
    while b - a > p:
        m = (a + b) / 2
        if f(a) * f(m) <= 0:
            b = m
        else:
            a = m
    return a, b

def f(x):
    return x**3 + x - 1

print(dichotomie(f, 0, 1, 0.01))
\end{lstlisting}
\end{minipage}\hfill
\begin{minipage}[t]{0.42\linewidth}
\begin{exemple}{}{}
Avec $f(x)=x^3+x-1$ sur $\intff{0}{1}$ et $p=0{,}01$, le programme affiche :
\begin{center}
\texttt{(0.6796875, 0.6875)}
\end{center}
Ainsi $0{,}679<\alpha<0{,}688$ après $7$ étapes
($\frac{1}{2^7}\approx 0{,}0078<0{,}01$), ce qui est cohérent avec l'encadrement
obtenu par balayage.
\end{exemple}
\end{minipage}

% =====================================================================
\section*{L'essentiel du chapitre}
\addcontentsline{toc}{section}{L'essentiel du chapitre}
% =====================================================================

\begin{center}
\renewcommand{\arraystretch}{1.6}
\begin{tabular}{|>{\raggedright\arraybackslash}m{4.2cm}|>{\raggedright\arraybackslash}m{11.6cm}|}
\hline
\rowcolor{bleuclair}\textbf{Notion} & \textbf{À retenir} \\ \hline
Continuité en $a$ & $\limite{x\to a}{f(x)}=f(a)$ ; la courbe se trace « sans lever le crayon ». \\ \hline
Fonctions usuelles & Polynômes, $\e^x$, $\sin$, $\cos$, $\lvert x\rvert$ continues sur $\R$ ; $\sqrt{x}$ sur $\intfo{0}{+\infty}$ ; rationnelles sur leurs intervalles de définition. \\ \hline
Opérations & Somme, produit, quotient (dénominateur non nul) et composée de fonctions continues sont continues. \\ \hline
Dérivabilité & Dérivable $\Longrightarrow$ continue ; la réciproque est fausse ($\lvert x\rvert$, $\sqrt{x}$ en $0$). \\ \hline
Suites & Si $u_n\to\ell$ et $f$ continue en $\ell$, alors $f(u_n)\to f(\ell)$ ; si $u_{n+1}=f(u_n)$ converge vers $\ell$, alors $f(\ell)=\ell$. \\ \hline
TVI & $f$ continue sur $\intff{a}{b}$, $k$ entre $f(a)$ et $f(b)$ : au moins une solution à $f(x)=k$. \\ \hline
Corollaire & Si de plus $f$ est strictement monotone : une \textbf{unique} solution. \\ \hline
\end{tabular}
\end{center}

\end{document}
